\documentclass[11pt,leqno]{article}
\headheight=13.6pt

% packages
\usepackage[alphabetic]{amsrefs}
\usepackage{physics}
% margin spacing
\usepackage[top=1in, bottom=1in, left=0.5in, right=0.5in]{geometry}
\usepackage{amsfonts, amsmath, amssymb, amsthm}
\usepackage{extarrows}
\usepackage[none]{hyphenat}
\usepackage{fancyhdr}
\usepackage{graphicx}
\graphicspath{{./images/}}
\usepackage{float}
\usepackage{color}
\newcommand{\sai}[1]{\textcolor{red}{#1}}
\usepackage{enumitem}
% \usepackage{mathrsfs}
\usepackage{hyperref}
\usepackage[noabbrev, capitalise]{cleveref}
\crefformat{equation}{equation~#2#1#3}
\crefformat{lemma}{\textrm{Lemma}~#2#1#3}
\usepackage{quiver}

% theorems
\theoremstyle{plain}
\newtheorem{lem}{Lemma}
\newtheorem{lemma}[lem]{Lemma}
\newtheorem{thm}[lem]{Theorem}
\newtheorem{theorem}[lem]{Theorem}
\newtheorem{prop}[lem]{Proposition}
\newtheorem{proposition}[lem]{Proposition}
\newtheorem{cor}[lem]{Corollary}
\newtheorem{corollary}[lem]{Corollary}
\newtheorem{conj}[lem]{Conjecture}
\newtheorem{fact}[lem]{Fact}
\newtheorem{form}[lem]{Formula}

\theoremstyle{definition}
\newtheorem{defn}[lem]{Definition}
\newtheorem{definition/}[lem]{Definition}
\newenvironment{definition}
  {\renewcommand{\qedsymbol}{\textdagger}%
   \pushQED{\qed}\begin{definition/}}
  {\popQED\end{definition/}}
\newtheorem{example}[lem]{Example}
\newtheorem{remark}[lem]{Remark}
\newtheorem{exercise}[lem]{Exercise}
\newtheorem{notation}[lem]{Notation}

\numberwithin{equation}{section}
\numberwithin{lem}{section}

% header/footer formatting
\pagestyle{fancy}
\fancyhead{}
\fancyfoot{}
\fancyhead[L]{Beilinson-Bernstein localization}
\fancyhead[C]{}
\fancyhead[R]{Sai Sivakumar}
\fancyfoot[R]{\thepage}
\renewcommand{\headrulewidth}{1pt}

% paragraph indentation/spacing
\setlength{\parindent}{0cm}
\setlength{\parskip}{10pt}
\renewcommand{\baselinestretch}{1.25}

% mathscript S
\usepackage[mathscr]{euscript}

% math operators
\DeclareMathOperator{\Stab}{Stab}
\DeclareMathOperator{\Ind}{Ind}
\let\O\relax
\DeclareMathOperator{\O}{\mathcal O}
\DeclareMathOperator{\D}{\mathcal D}
\newcommand{\catname}[1]{{\normalfont\mathbf{#1}}}
\newcommand{\Mod}{\text{-}\catname{mod}}
\DeclareMathOperator{\Lie}{Lie}

% bracket notation for inner product
\usepackage{mathtools}

\DeclarePairedDelimiterX{\abr}[1]{\langle}{\rangle}{#1}

% smileys frownies
\usepackage{wasysym}
\newcommand{\smallhappy}{\raisebox{-.14em}{\smiley}}
\newcommand{\happy}{\raisebox{-.24em}{\resizebox{1.2em}{!}{\smiley}}}
\newcommand{\smallsad}{\raisebox{-.14em}{\frownie}}
\newcommand{\sad}{\raisebox{-.24em}{\resizebox{1.2em}{!}{\frownie}}}
\DeclareMathOperator{\mathhappy}{\!\happy\!}
\DeclareMathOperator{\smallmathhappy}{\!\smallhappy\!}
\DeclareMathOperator{\mathsad}{\!\sad\!}
\DeclareMathOperator{\smallmathsad}{\!\smallsad\!}

% set page count index to begin from 1
\setcounter{page}{1}

% array column and row separation
\arraycolsep = 2pt
\renewcommand{\arraystretch}{.6}

% footnote style
\renewcommand{\thefootnote}{[\arabic{footnote}]}

\begin{document}

We state the Beilinson-Bernstein localization theorem and provide a couple examples. These notes closely follow (i.e. is a rough translation/summary of) \textit{Localisation de $\mathfrak{g}$-modules} by Beilinson and Bernstein and borrow examples from \textit{Four examples of Beilinson-Bernstein localization} by Romanov.

\section{$\mathcal D$-affine varieties and sheaves of twisted differential operators}
Let $X$ be a smooth variety over an algebraically closed field $k$ of characteristic $0$. We write $\D_X$ for the sheaf of differential operators on $X$. Denote by $\D_X\Mod$ the category of quasicoherent left $\D_X$-modules on $X$. By a quasicoherent $\D_X$-module we mean a $\D_X$-module which is quasicoherent as an $\O_X$-module.
\begin{definition}
  The variety $X$ is \textit{$\D$-affine} if every $\D_X$-module $\mathcal M$ is generated over $\D_X$ by its global sections and $H^i(X,\mathcal M) = 0$ for $i>0$.
\end{definition}
\begin{proposition}
  If $X$ is $\D_X$-affine, the global sections functor $\Gamma\colon \mathcal M\mapsto \Gamma(X,\mathcal M)$ defines an equivalence of $\D_X\Mod$ with the category of left $\Gamma(X,\D_X)$-modules.
\end{proposition}
\sai{TODO} proof

It is clear that every affine variety is $\D$-affine, but the converse is false:
\begin{theorem}
  The space $\mathbb P^n$ is $\D$-affine.
\end{theorem}
This result follows from the more general result \cref{thm: flag d-affinity} that all flag varieties are $\D$-affine.

We will need sheaves of rings on $X$ which slightly generalize the sheaf of differential operators $\D_X$. Consider the category of pairs $(\mathcal A,f)$, where $\mathcal A$ is a sheaf of $k$-algebras on $X$ and $f\colon \O_X\to\mathcal A$ is a morphism of sheaves of $k$-algebras. There is the natural inclusion $i\colon \O_X\to\D_X$, so it follows that the pair $(\D_X,i)$ belongs to the aforementioned category.
\begin{definition}
  The pair $(\mathcal A,f)$ is called a sheaf of \textit{twisted differential operators} if $(\mathcal A,f)$ is locally isomorphic to the pair $(\D_X,i)$. When it is clear from context we will refer to the pair $(\mathcal A,f)$ by $\mathcal A$.
\end{definition}
\begin{remark}
  The automorphisms of $\D_X$ are determined by their restriction to the first-order differential operators in $\D_X$.\footnote{\sai{TODO} The sheaf $\D_X$ is generated by the tangent sheaf $\mathcal T_X$ of $X$ over $\O_X$.} Such automorphisms are given by $v\mapsto v-i(\abr{\alpha,v})$ for $\alpha$ a closed $1$-form on $X$ (where $v$ is a first-order differential operator). Sheaves of twisted differential operators correspond bijectively to $\Omega_{X,\text{cl}}^1$-torsors, where $\Omega_{X,\text{cl}}^1$ is the sheaf of closed $1$-forms on $X$.
\end{remark}
\begin{example}
  Let $\mathcal L$ be an invertible $\O_X$-module. Then the sheaf of algebras $\D_{\mathcal L}$ of differential operators from $\mathcal L$ to itself constitutes a sheaf of twisted differential operators. The corresponding $\Omega_{X,\text{cl}}^1$-torsor is the logarithmic differential of the $\O_X^\ast$-torsor corresponding to $\mathcal L$.\footnote{\sai{TODO} See .} 
\end{example}
For a pair $(A,i)$ as before, denote by $\mathcal A\Mod$ the category of quasicoherent (as $\O_X$-modules) $\mathcal A$-modules; similarly we say that $X$ is $\mathcal A$-affine if every $\mathcal M\in \mathcal A\Mod$ is generated by its global sections and $H^i(X,\mathcal M) = 0$ for $i>0$.

\section{The localization theorem}
Let $G$ be a connected reductive group over $k$, $\mathfrak g = \Lie(G)$ its Lie algebra, $U(\mathfrak g)$ the enveloping algebra of $\mathfrak g$, and $Z(\mathfrak g)$ the center of $U(\mathfrak g)$. Let $X$ be the complete flag variety of $G$; that is, $X$ is the variety $G/B$ which parameterizes the Borel subgroups of $G$.\footnote{\sai{TODO} explain.} Let $\mathcal T_X$ be the tangent sheaf of $X$, $\alpha\colon \mathfrak g\to\mathcal T_X$ the morphism of sheaves of Lie algebras induced by the (\sai{TODO} left) action of $G$ on $X$. Let $B_x$ be the Borel subgroup of $G$ corresponding to $x\in X$, $N_x$ its unipotent radical, and $H=B_x/N_x$ a Cartan subgroup (the Cartan subgroups $B_x/N_x$ for each $x\in X$ are canonically isomorphic). Let $\mathfrak b_x = \Lie(B_x)$, $\mathfrak n_x = \Lie(N_x)$, $\mathfrak h = \Lie(H)$.
\end{document}